\documentclass[../../../include/open-logic-section]{subfiles}

\begin{document}

\olfileid[pt]{sth}{card-arithmetic}{fix}
\olsection{Pontos Fixos de $\aleph$}

No \olref[spine][]{chap}, sugerimos que o Esquema de Substituição carece de
justificação, pois força a hierarquia a ser bastante alta.
Tendo feito alguma aritmética cardinal, podemos dar uma pequena
ilustração da altura da hierarquia. 

Evidentemente $0 < \aleph_0$, e $1 < \aleph_1$, e $2 < \aleph_2$\ldots
e, de fato, a diferença de tamanho só fica \emph{maior} a cada
passo. Assim, é tentador conjecturar que $\kappa< \aleph_\kappa$
para todo o ordinal $\kappa$. 

Mas esta conjectura é \emph{falsa}, em $\ZFC$. De fato, podemos
provar que existem \emph{pontos fixos de $\aleph$}, isto é,
cardinais $\kappa$ tais que $\kappa=\aleph_\kappa$. 

\begin{prop}\ollabel{alephfixed}
Existe um ponto fixo de $\aleph$.
\end{prop}

\begin{proof}
Por recursão, defina-se:
\begin{align*}
	\kappa_0 &= 0\\
	\kappa_{n+1} &= \aleph_{\kappa_n}\\
	\kappa&= \bigcup_{n < \omega}\kappa_n
\end{align*}
Ora, $\kappa$ é um cardinal por
\olref[cardinals][classing]{unioncardinalscardinal}. Mas então:
\[
	\kappa= \bigcup_{n < \omega} \kappa_{n+1} = 
	\bigcup_{n < \omega}\aleph_{\kappa_n} = 
	\bigcup_{\alpha < \kappa}\aleph_\alpha = \aleph_\kappa
\]
%	Por construção, $\kappa$ é o menor cardinal maior do que cada
%	$\kappa_n$. Logo $\aleph_\kappa$ é o menor cardinal maior do que
%	cada $\aleph_{\kappa_n}$, e portanto maior do que cada
%	$\kappa_{n+1}$. Mas igualmente $\kappa$ é o menor cardinal maior
%	do que cada $\kappa_{n+1} = \aleph_{\kappa_n}$. Logo $\kappa=
%	\aleph_\kappa$.
\end{proof}

Boolos escreveu um artigo sobre precisamente o ponto fixo de $\aleph$ que
acabamos de construir. Depois de notar a existência de $\kappa$, no início
do artigo, disse:
\begin{quote}
	[$\kappa$ é] um número \emph{bastante grande}, aos olhos de quem
	não teve contato prévio com a teoria dos conjuntos, tão grande, me parece,
	que põe em causa a verdade de qualquer teoria que, entre as suas
	afirmações, inclua a de que existem pelo menos $\kappa$ objetos.
	\citep[p.~257]{Boolos2000}
\end{quote}
E concluiu o artigo perguntando:
\begin{quote}
	[será que] suspeitamos que, seja qual for o começo da história,
	quando chegamos tão longe as rodas giram no vazio e já não estamos a ouvir
	uma descrição de coisa alguma que seja o caso?
	\citep[p.~268]{Boolos2000}
\end{quote}
Se, de fato, ultrapassamos ``qualquer coisa que seja o caso'', então devemos
apontar o dedo da culpa diretamente ao Esquema de Substituição. Pois é este
axioma que permite a nossa prova funcionar. Nesse caso, supõe-se,
Boolos teria de rever a afirmação que fizera, décadas antes,
de que a Substituição não tem consequências ``indesejáveis'' (ver
\olref[replacement][extrinsic]{sec}).

Mas será a existência de $\kappa$ assim tão má? Pode ajudar, aqui,
considerar o \emph{paradoxo de Tristram Shandy} de Russell. Tristram Shandy
documenta a sua vida no diário, mas demora um ano a registrar um
único dia. A cada ano que passa, Tristram fica cada vez mais
atrasado: passado um ano, só registrou um dia, e viveu 364 dias por registrar; passados dois anos, só registrou
dois dias, e viveu 728 dias por registrar; passados três anos, só
registrou três dias, e viveu 1092 dias por registrar
\dots\footnote{Ignorando anos bissextos.} Ainda assim, se Tristram
é \emph{imortal}, conseguirá registrar todos os dias, pois
registrará o $n$-ésimo dia no $n$-ésimo ano da sua vida. E assim, ``no
fim dos tempos'', Tristram terá um diário completo. 

Porque é que isto difere tanto da ideia de que $\alpha$ é
menor do que $\aleph_\alpha$---e, de fato, cada vez mais, desesperadamente
menor---até $\kappa$, onde nos ``apanhamos'' e $\kappa
= \aleph_\kappa$?

Deixando isto de lado, e supondo que aceitamos $\ZFC$, fechemos com um
pouco mais de diversão sobre construções de pontos fixos. Os três resultados seguintes estabelecem, intuitivamente, que existe um ponto (não trivial) em
que a hierarquia é tão larga como alta:

\begin{prop}\ollabel{bethfixed}
Existe um ponto fixo de $\beth$, isto é, um $\kappa$ tal que $\kappa=
\beth_\kappa$.
\end{prop}

\begin{proof}
Como na \olref{alephfixed}, usando ``$\beth$'' no lugar de ``$\aleph$''. 
\end{proof}

\begin{prop}\ollabel{stagesize}
$\card{V_{\omega+\alpha}} = \beth_{\alpha}$. Se $\omega \ordtimes
\omega \leq \alpha$, então $\card{V_\alpha} = \beth_\alpha$.
\end{prop}

\begin{proof}
A primeira afirmação vale por uma indução transfinita simples. A segunda
segue-se, pois se $\omega \ordtimes \omega \leq \alpha$ então
$\omega + \alpha = \alpha$. Para o estabelecer, usamos fatos sobre
aritmética ordinal de \olref[ord-arithmetic][]{chap}. Note-se primeiro
que $\omega \ordtimes \omega = \omega \ordtimes (1 \ordplus \omega) =
(\omega  \ordtimes 1) \ordplus (\omega\ordtimes\omega) = \omega
\ordplus (\omega \ordtimes \omega)$. Ora, se $\omega \ordtimes \omega
\leq \alpha$, isto é, $\alpha = (\omega\ordtimes\omega) \ordplus \beta$
para algum $\beta$, então $\omega \ordplus \alpha = \omega \ordplus
((\omega \ordtimes \omega) \ordplus \beta) = (\omega \ordplus (\omega
\ordtimes \omega)) \ordplus \beta = (\omega \ordtimes \omega) \ordplus
\beta = \alpha$. 
\end{proof}

\begin{cor}
Existe um $\kappa$ tal que $\card{V_\kappa} = \kappa$.
\end{cor}

\begin{proof}
Seja $\kappa$ um ponto fixo de $\beth$, dado por \olref{bethfixed}.
Claramente $\omega \ordtimes \omega < \kappa$. Logo $\card{V_\kappa} =
\beth_\kappa= \kappa$ por \olref{stagesize}.
\end{proof}

Há tantos estágios abaixo de $V_\kappa$ quantos !!{element}s
de $V_\kappa$. Intuitivamente, então, $V_\kappa$ é tão largo como alto.
Isto é muito à maneira de Tristram Shandy: passamos de um estágio ao seguinte
tomando \emph{conjuntos potência}, tornando a hierarquia \emph{muito}
maior a cada passo. Mas ``no fim'', isto é, no estágio $\kappa$,
a largura da hierarquia alcança a altura. 

Alguém pode perguntar: \emph{Com que frequência a largura da hierarquia coincide com a
altura?} A resposta é: \emph{Tantas vezes quantos os ordinais.} Mas
isto precisa de uma pequena explicação. 

Definimos um termo $\tau$ do seguinte modo. Para qualquer $A$, seja:
\begin{align*}
	\tau_0(A) & \defis \card{A}\\
	\tau_{n+1}(A) & \defis \beth_{\tau_n(A)}\\
	\tau(A) & \defis \bigcup_{n < \omega}\tau_n(A)
\intertext{Como na \olref{bethfixed}, $\tau(A)$ é um
ponto fixo de $\beth$ para qualquer $A$, e trivialmente $\card{A} < \tau(A)$.
Considere-se então esta definição recursiva:}
	W_0 &\defis 0\\
	W_{\alpha + 1} & \defis \tau(W_\alpha)\\
	W_\alpha & \defis \bigcup_{\beta < \alpha} W_\beta
	\text{, quando $\alpha$ é limite}
\end{align*}
A construção está definida para todos os ordinais. Intuitivamente, então,
$W$ é ``!!a{injection}'' dos ordinais nos pontos fixos de $\beth$.
E, exatamente como antes, 
$V_{W_\alpha}$ é tão largo como alto, para qualquer $\alpha$.

\end{document}
